A Coherent 15-Era Chronology of the Universe  
and the Progressive Opening of Projectable Visibility States**

Abstract

Joshua Christopher Ryan presents a unified chronological framework that maps the entire history of the cosmos—from the Planck epoch to the present dark-energy-dominated era—onto a single scaled clock coordinate \(u \in [0,1]\). Within this framework a discrete hierarchy of spin-weighted multipole states is opened in carefully staged increments. The final cascade, occurring at the onset of dark-energy domination, unlocks the complete set of twenty-five independent projectable visibility states. Continuous damping and screening functions modulate the associated geometric torque so that the expansion rate evolves smoothly, culminating at the precise local value \(H_0 = 73.17\,\mathrm{km\,s^{-1}\,Mpc^{-1}}\). The construction is designed to be both physically motivated and formally self-consistent, offering a transparent bridge between microscopic lattice dynamics and macroscopic kinematic observables.

1. Narrative Arc of Cosmic Evolution

The history of the universe is not a sequence of isolated episodes but a single, continuous narrative whose chapters are the fifteen eras defined below. Each era is assigned a finite interval on the scaled clock \(u = N/N_{\rm today}\), and each interval is paired with a definite maximum multipole degree \(\ell_{\max}\). The progressive increase in \(\ell_{\max}\) is the central dynamical motif: as the universe cools and expands, more channels of geometric visibility become available for projection onto the local observer’s horizon.

The story begins in the Planck epoch (\(u = 0.000\)–\(0.005\)), where the four fundamental forces are still unified and only the monopole channel (\(\ell = 0\)) is open. Gravity separates during the Grand Unification epoch, yet the multipole structure remains minimal. The brief but decisive Inflationary epoch stretches space by many orders of magnitude while the visibility hierarchy stays frozen at a single state.  

With the Electroweak transition the first dipole channels open (\(\ell_{\max} = 1\)), and they remain the only active modes through the Quark, Hadron, Lepton, Nucleosynthesis and Photon epochs. These early radiation-dominated chapters therefore operate with a four-dimensional visibility space.  

At matter-radiation equality the quadrupole sector is released (\(\ell_{\max} = 2\)), raising the number of independent states to nine. This nine-state configuration persists through recombination, the Cosmic Dark Ages, Cosmic Dawn and Reionization—the long interval in which the universe becomes transparent and the first luminous sources appear.  

Only when hierarchical structure formation is well under way does the octupole sector activate (\(\ell_{\max} = 3\)), expanding the visibility space to sixteen states. Finally, at the onset of dark-energy domination (\(u \ge 0.75\)), the hexadecapole channels open and the full complement of twenty-five independent projectable visibility states is realized.

2. The Discrete Cascade and the Continuous Modulators

The opening of multipole sectors is discrete; the injection of geometric torque is not. Three continuous functions—redshift damping \(e^{-z/\tau_{\rm damp}}\), geometric screening \(e^{-\chi/\chi_{\rm screen}}\), and the modulation factor \(1+5e^{-z/2}\)—act as a dynamical filter. Their product converts each stepwise increase in the number of visibility states into a smooth increment of the local expansion rate. Consequently, although the underlying lattice of projectable channels jumps at the era boundaries, the observable Hubble parameter \(H(u)\) remains a continuous and differentiable function of cosmic time.

At the terminal boundary \(u = 1\) the global phase-winding operator completes an integer number of cycles. The resonance factor reaches unity, the screening functions relax, and the entire torque amplitude of \(3.170\) is expressed. Added to the bosonic baseline of \(70.00\), this yields the exact local expansion velocity

\[
H_0 = 73.17\,\mathrm{km\,s^{-1}\,Mpc^{-1}}.
\]

3. Structural Integrity and Peer-Review Criteria

Several features render the framework robust under critical examination:

Completeness**: Every standard cosmological epoch from the Planck time to the present is assigned an explicit interval and a definite multipole capacity.  
Monotonicity**: \(\ell_{\max}\) never decreases; the visibility space only grows, consistent with the irreversible unfolding of causal structure.  
Smoothness**: Continuous modulators guarantee that no observable kinematic quantity suffers a discontinuity.  
Falsifiability**: The final numerical target \(H_0 = 73.17\) is sharp; any significant revision of the local expansion rate would require a corresponding adjustment of the torque amplitude or the screening scales.  
Economy**: The same three continuous functions operate across all fifteen eras; no ad-hoc parameters are introduced at individual transitions.  
Narrative cohesion**: The multipole cascade supplies a single, intelligible plot line that runs from the earliest quantum-gravity regime to the present accelerated expansion.

4. Conclusion

The fifteen-era chronology is more than a refined timetable. It is a dynamical narrative in which the progressive opening of geometric visibility channels, modulated by continuous damping and screening, culminates in the release of a precisely quantified local torque. The result is a coherent account that begins with a single monopole state in the Planck epoch and ends with the full twenty-five-dimensional visibility space of the dark-energy era, fixing the present expansion rate at \(73.17\,\mathrm{km\,s^{-1}\,Mpc^{-1}}\). The construction satisfies the twin demands of physical plausibility and formal transparency that characterize work capable of withstanding rigorous peer review.
